\section*{Chapitre \ref{ch:b1:mammal-organization} --- Organisation fonctionnelle d'un mammifère}

\begin{solution}{exo:b1:mammal-organization:1}
Vertébré~: symétrie bilatérale, \omterm{def:b1:mammal-organization:bodyplan}{céphalisation}, axe vertébral dorsal
entourant la moelle épinière, \omterm{def:b1:mammal-organization:bodyplan}{cœlome} ventral contenant les viscères,
quatre membres. Mammalien~: le \omterm{def:b1:mammal-organization:bodyplan}{diaphragme} divisant le \omterm{def:b1:mammal-organization:bodyplan}{cœlome} en thorax
et abdomen, les poils, le lait, la température corporelle régulée.
\end{solution}

\begin{solution}{exo:b1:mammal-organization:2}
Nutrition~: digestif (estomac), respiratoire (poumon), circulatoire
(cœur), excréteur (rein). Relation~: nerveux (encéphale), \omterm{def:b1:mammal-organization:organ}{organes} des
sens (œil), endocrinien (thyroïde), squelette et muscles (fémur,
biceps), peau. Reproduction~: génital (ovaire, testicule), glandes
mammaires.
\end{solution}

\begin{solution}{exo:b1:mammal-organization:3}
Eau $0.6\times 60 = \qty{36}{L}$~; intracellulaire \qty{24}{L}~;
extracellulaire \qty{12}{L}, dont interstitiel \qty{9}{L} et \omterm{def:b1:mammal-organization:compartments}{plasma}
\qty{3}{L}.
\end{solution}

\begin{solution}{exo:b1:mammal-organization:4}
$1.1/5 = 22\%$. \qty{1.1}{L/min} font \qty{1580}{L} par jour, soit 290
fois la volémie.
\end{solution}

\begin{solution}{exo:b1:mammal-organization:5}
Inuline restante \qty{1.7}{g}~; $V = 1.7/0.12 = \qty{14}{L}$ pour
l'extracellulaire~; intracellulaire $42 - 14 = \qty{28}{L}$.
\end{solution}

\begin{solution}{exo:b1:mammal-organization:6}
La paroi capillaire laisse librement passer l'eau et les petits ions,
qui s'équilibrent donc, mais retient les protéines. Les protéines
plasmatiques attirent l'eau dans les vaisseaux par osmose~; sans elles
l'eau fuirait vers l'espace interstitiel, qui gonflerait (œdème), et le
volume plasmatique diminuerait.
\end{solution}

\begin{solution}{exo:b1:mammal-organization:7}
En parallèle, chaque \omterm{def:b1:mammal-organization:organ}{organe} reçoit du sang à la pleine teneur artérielle
en oxygène et à la pleine pression, et son débit peut être réglé
indépendamment. En série, l'encéphale recevrait un sang déjà appauvri
par les muscles, à plus basse pression, et son approvisionnement
chuterait précisément à l'effort.
\end{solution}

\begin{solution}{exo:b1:mammal-organization:8}
$hS\Delta T = 80 - 20 = \qty{60}{W}$~; $h = 60/(1.8\times 13) =
\qty{2.6}{W.m^{-2}.K^{-1}}$. Dans un air à \qty{33}{\celsius} la perte
non évaporative s'annule~; la peau doit s'échauffer au-dessus de l'air
(la vasodilatation la porte vers \qty{36}{\celsius}) et la sudation doit
emporter le reste.
\end{solution}

\begin{solution}{exo:b1:mammal-organization:9}
Chaleur à évacuer $800 - 200 - 150 = \qty{450}{W}$, soit
\qty{1620}{kJ/h}, \qty{675}{g} de sueur par heure. Sans sudation,
$450/(70\times 3500) =
\qty{1.8e-3}{K/s}$, soit environ \qty{0.11}{K/min}~: \qty{3}{K} en une
demi-heure.
\end{solution}

\begin{solution}{exo:b1:mammal-organization:10}
Les petits corps perdent la chaleur le plus vite par gramme (grand
$S/V$) et ont besoin d'une production continue, élevée et fiable~: la
graisse brune produit de la chaleur chimiquement, à débit constant et
sans mouvement, tandis que le frisson exige des muscles développés (que
les nouveau-nés n'ont pas), gêne le déplacement et la prise alimentaire,
et est intermittent. La graisse brune est de plus placée près des gros
vaisseaux et réchauffe le sang directement.
\end{solution}

\begin{solution}{exo:b1:mammal-organization:11}
L'hypothalamus réchauffé lit «~trop chaud~»~: le chien halète et dilate
ses vaisseaux cutanés malgré le froid, et sa température centrale chute
d'un ou deux degrés en une heure. Sans hypothalamus, il n'y a plus de
\omterm{def:b1:organism-environment:homeostasis}{valeur de consigne}~: le chien ne frissonne pas au froid et ne halète pas
à la chaleur, et sa température dérive vers celle de la pièce dans les
deux cas.
\end{solution}

\begin{solution}{exo:b1:mammal-organization:12}
Le \omterm{def:b1:mammal-organization:compartments}{liquide extracellulaire} (\omterm{def:b1:mammal-organization:compartments}{liquide interstitiel} et \omterm{def:b1:mammal-organization:compartments}{plasma}), de
composition riche en sodium, de température et de volume presque
constants, baigne chaque cellule comme l'eau de mer baignait les
premiers animaux. Sa constance est maintenue par des boucles de
\omterm{def:b1:organism-environment:homeostasis}{rétroaction négative} (le rein pour le volume et les ions, le poumon pour
les gaz, le foie et le pancréas pour le glucose, l'hypothalamus pour la
température). Ses «~rivages~» sont les \omterm{prop:b1:organism-environment:surfaces}{surfaces d'échange} --- intestin,
poumon, rein, peau --- où il rencontre l'extérieur, et la circulation
est le courant qui le brasse.
\end{solution}

\begin{solution}{pb:b1:mammal-organization:1}
\textbf{1.} $B = 3.4\times 2^{0.75} = \qty{5.7}{W}$.
\textbf{2.} Terrier~: $5.7\times 12\times 3600 = \qty{246}{kJ}$~;
recherche de nourriture~: $17.1\times 6\times 3600 = \qty{369}{kJ}$~;
dehors~: $11.4\times 6\times 3600 = \qty{246}{kJ}$.
\textbf{3.} \qty{861}{kJ}, soit $861/493 = 1.75$ fois le \omterm{def:b1:organism-environment:allometry}{métabolisme
basal} ($B$ sur \qty{24}{h} $= \qty{493}{kJ}$).
\textbf{4.} $0.25\times 18\times 0.65 = \qty{2.9}{kJ/g}$.
\textbf{5.} $861/2.9 = \qty{295}{g}$ d'herbe fraîche, \qty{74}{g} de
matière sèche.
\textbf{6.} $295/2000 = 15\%$ de sa masse.
\textbf{7.} $861/20 = \qty{43}{L}$ d'$\mathrm{O_2}$~; \qty{30}{mL/min}.
\textbf{8.} $0.9\times 43 = \qty{39}{L}$ de $\mathrm{CO_2}$, soit $39/22.4\times 44 = \qty{76}{g}$.
\textbf{9.} $0.75\times 295 = \qty{221}{g}$.
\textbf{10.} Matière sèche digérée $0.65\times 74 = \qty{48}{g}$~; eau
$0.6\times 48 = \qty{29}{g}$.
\textbf{11.} Non digérée \qty{26}{g}~; fèces $26/0.6 = \qty{43}{g}$,
emportant \qty{17}{g} d'eau.
\textbf{12.} $(44 - 5)\times 0.86 = \qty{34}{g}$.
\textbf{13.} Entrées~: $221 + 29 = \qty{250}{g}$. Sorties~: $17 + 130 + 34 + 20
= \qty{201}{g}$. Excédent \qty{49}{g}.
\textbf{14.} Non~: l'herbe apporte plus qu'il ne perd~; l'excédent part
en urine supplémentaire (environ \qty{180}{mL} en tout).
\textbf{15.} Même matière sèche \qty{74}{g} dans $\qty{148}{g}$
d'herbe~: eau alimentaire \qty{74}{g}~; entrées \qty{103}{g} contre
sorties \qty{201}{g}~: un déficit de \qty{100}{g} par jour, qu'il doit
boire ou prélever sur son corps.
\textbf{16.} Eau $0.6\times 2.0 = \qty{1.2}{L}$~; extracellulaire
\qty{0.4}{L}~; \omterm{def:b1:mammal-organization:compartments}{plasma} \qty{0.1}{L}.
\textbf{17.} $1.5/15 = \qty{0.10}{L}$~: en accord.
\textbf{18.} $0.10/0.60 = \qty{0.167}{L}$, \qty{8}{\%} de la masse corporelle.
\textbf{19.} $5\times(2/70)^{0.75} = \qty{0.35}{L/min}$~; temps de
circuit $0.167/0.35 = \qty{0.48}{min}$, soit environ \qty{29}{s}.
\textbf{20.} $350/200 = \qty{1.75}{mL}$ par battement.
\textbf{21.} $0.1\times 2^{2/3} = \qty{0.16}{m^2}$.
\textbf{22.} $2.5\times 0.16\times 34 = \qty{13.6}{W}$, contre
$2B = \qty{11.4}{W}$ supposés~: le lapin doit produire un peu plus, ou
réduire ses pertes par la posture et l'abri.
\textbf{23.} Ouvertes~: $10\times 0.01\times 34 = \qty{3.4}{W}$, soit
encore un quart de la perte par le pelage~; fermées~: $10\times 0.01\times 3 =
\qty{0.3}{W}$. Il les ferme.
\textbf{24.} Pelage~: $2.5\times 0.16\times 9 = \qty{3.6}{W}$~;
oreilles ouvertes~: $10\times 0.01\times 9 = \qty{0.9}{W}$~; total
\qty{4.5}{W} contre \qty{5.7}{W} produits. Les \qty{1.2}{W} restants
doivent partir par évaporation~: le lapin halète (\qty{1.2}{W}
représentent \qty{1.8}{g} d'eau par heure), s'aplatit sur un sol frais
et reste à l'ombre.
\textbf{25.} Dépense énergétique quotidienne \qty{0.86}{MJ}, soit $1.75$
fois son \omterm{def:b1:organism-environment:allometry}{métabolisme basal}~; renouvellement d'eau \qty{250}{g}, soit
\qty{12.5}{\%} de sa masse.
\end{solution}
